\section{El mercado de Agent y la flota de las grandes empresas}

El riesgo de la complejidad, tratado antes, se refiere al interior de Manus; esta sección mira el mercado externo. En la segunda mitad de la entrevista, Peak evaluó a OpenAI, Anthropic, Google Gemini, xAI, Meta, Thinking Machines, la nueva empresa de Ilya, entre otras. Sus observaciones tienen un punto en común: la competencia entre modelos no tiene un punto final, la competencia entre productos tampoco tiene un final simple, pero cada empresa se diferenciará a partir de sus propias fortalezas.

\begin{figure}[H]
\centering
\includegraphics[width=0.94\textwidth]{agent-market-map.png}
\caption{Mapa del mercado de Agent: Coding, Agent verticales, Agent para el consumidor final, ecosistema complementario y empresas de modelos.}
\end{figure}

\begin{knowledgebox}{Lectura de la figura: el mercado de Agent no es una sola pista de competencia}
Los Coding Agent, los Agent verticales, los Agent para el consumidor final, el ecosistema de pagos y permisos y las empresas de modelos avanzan todos a la vez. Manus se ubica más cerca del punto entre el consumidor final y las herramientas de productividad: necesita capacidad del modelo, pero también un lugar en la mente del usuario y flujos de trabajo reales.
\end{knowledgebox}

\subsection{La IA ya no compite solo por el tiempo del usuario}

Peak plantea un juicio muy interesante: el internet móvil terminó compitiendo por el tiempo total del usuario, mientras que un Agent puede crear más valor al tiempo que reduce el tiempo de interacción directa. El Agent trabaja de forma continua en segundo plano y el usuario no necesita vigilarlo todo el tiempo. Por eso, el estado estable de la era de la IA no necesariamente estará determinado, como en el internet móvil, por la «cantidad total de atención».

\begin{knowledgebox}{En qué se distingue el Agent del internet móvil}
Los productos del internet móvil suelen disputarse el tiempo del usuario; los productos de Agent quizá se disputen la autorización del usuario, el punto de entrada de las tareas, la ejecución en segundo plano y la confiabilidad de los resultados. No necesariamente requieren que el usuario esté conectado todo el tiempo.
\end{knowledgebox}

\subsection{Resumen de la sección}

El mercado de Agent se diferenciará en varios niveles. La oportunidad de Manus no depende solo de la capacidad del modelo, sino de si logra convertirse en el principal punto de entrada a la productividad del usuario sin dejar de ser simple y distintivo.
